\markdownRendererDocumentBegin
\markdownRendererSectionBegin
\markdownRendererHeadingOne{Introduction}\markdownRendererInterblockSeparator
{}Test-time computation has become an explicit design variable in modern reasoning systems. Systems allocate more tokens, samples, search nodes, verifier calls or iterative refinement to difficult instances. Recent work develops bandit, constrained-policy and learned adaptive allocation strategies, reinforcing a general lesson: uniform inference budgets waste computation when item difficulty is heterogeneous.\markdownRendererInterblockSeparator
{}But difficulty itself has more than one source. Some tasks are difficult because the relevant structure is poorly exposed in the current representation. Others are difficult because substantial search or reasoning remains even after the right structure is visible. If a system spends all marginal budget on reasoning in an access-limited task, it reasons harder over the wrong state. If it spends all marginal budget on state construction in a reasoning-limited task, it repeatedly reorganizes information that was already accessible.\markdownRendererInterblockSeparator
{}P12 asks:\markdownRendererInterblockSeparator
{}\markdownRendererBlockQuoteBegin
\markdownRendererStrongEmphasis{Under one matched total budget, when should a system spend computation changing state, when should it spend computation reasoning over state, and can a prospective policy learn or exploit the difference?}
\markdownRendererBlockQuoteEnd \markdownRendererInterblockSeparator
{}The system is modeled as\markdownRendererInterblockSeparator
{}\markdownRendererCodeSpan{raw/current state -> optional state construction -> downstream reasoning/search -> verified outcome}.\markdownRendererInterblockSeparator
{}The paper makes three contributions.\markdownRendererInterblockSeparator
{}\markdownRendererOlBeginTight
\markdownRendererOlItemWithNumber{1}\markdownRendererStrongEmphasis{Two-axis inference formulation.} State construction and downstream reasoning are symmetric budgeted actions rather than free preprocessing plus paid reasoning.\markdownRendererOlItemEnd 
\markdownRendererOlItemWithNumber{2}\markdownRendererStrongEmphasis{A strict comparator contract.} The joint policy must beat both adaptive-state-only and adaptive-reasoning-only policies at identical total resources, not merely a fixed baseline.\markdownRendererOlItemEnd 
\markdownRendererOlItemWithNumber{3}\markdownRendererStrongEmphasis{Protected held-out evidence.} In 16 held-out resource families, a frozen joint allocator produces a +0.3347 average advantage over the stronger one-axis adaptive comparator, with positive gain in every family.\markdownRendererOlItemEnd 
\markdownRendererOlEndTight \markdownRendererInterblockSeparator
{}The result is intentionally controlled. Its purpose is to establish that joint state–reasoning allocation has a falsifiable value beyond generic adaptive test-time compute, and to define the matched-budget contract a real-system paper must satisfy.
\markdownRendererSectionEnd \markdownRendererDocumentEnd